\section{遇到的困難與心得}

% 草稿：以下依開發紀錄整理，請改寫成自己的感受

\subsection{從 GUI 框架到遊戲架構}

HW2 一開始沿用 HW1 的程式，但 HW1 為了繪圖程式寫的 Element Tree、事件分派與多視窗，放到一個持續運作的 3D 場景裡並不合用：機器人需要每個 tick 推進的模擬，而不是等待事件才重畫。
因此我只保留字型、數學、平台層等基礎模組，其餘改寫成以固定 tick 推進的遊戲架構，並規定 \lstinline|world| 不能依賴任何繪圖或輸入的程式。
這個限制一開始讓很多事情變得麻煩，例如相機要的眼睛位置不能從繪圖端取得，只好把骨架與正向運動學放進 \lstinline|world|；但也因為如此，碰撞、姿勢與相機都使用同一份骨架，不會出現畫面與模擬不一致的情形。

\subsection{碰撞處理}

碰撞是除錯最久的部分。單看偵測公式並不難，困難在於各種邊界情況：
站在方塊頂端時，表面相切的情況沒有算成接觸，機器人會在著地與懸空之間不停切換；
走下坡時每一步都略微離地，被當成在空中而無法控制；
姿勢改變後碰撞箱還是站立時的大小，坐下時腳會穿進障礙物。
這些問題分別以「相切視為深度 0 的接觸」、「貼地距離」以及「依姿勢內插並隨朝向旋轉的碰撞箱」解決。
我也體會到，物理模擬中「看起來對」需要處理的細節，遠比公式本身多。

\subsection{從階層模型到骨架}

一開始我以為階層模型只要用 Push/Pop 與局部座標系把每個部位畫出來就好；實際做了才發現，要讓機器人動起來，關節角度必須是資料而不是寫死在繪圖指令裡。
最後我把機器人拆成骨架（關節的樹與角度範圍）、姿勢（每個關節的角度）與外觀（每根骨頭的網格），姿勢之間可以內插，走路動畫只是在姿勢上疊加角度，關節限制也能在同一個地方統一處理。

\subsection{心得}

這次作業規定的內容其實只需要一塊平地和幾個障礙物，但我花了不少時間在地形生成上，調整雜訊、生態系與效能。
雖然超出作業範圍，卻讓我對 fBm、網格快取等技巧有了實際的體會，也讓機器人走在高低起伏的地面上時，碰撞與貼地的處理更容易測出問題。
回頭看，最有收穫的是把「意圖、模擬、繪製」分開的設計：鍵盤只負責產生意圖，速度與角度的限制由機器人自己決定，畫面只是把狀態畫出來。
功能越加越多時，這樣的分工讓我能放心修改其中一部分，而不必擔心影響其他地方。
